% Stil vs. form: overflatestatistikk (maskina sitt domene) mot målbart formrom
% (morfometri). For kapittelet «Den siste metoden». Reint TikZ.
\begin{figure}[t]
\centering
\begin{tikzpicture}[font=\footnotesize]
% ---- Venstre: stilklassifisering (overflate) ----
\node[anchor=north west, font=\small\scshape] at (0,3.3) {stil: overflatestatistikk};
\foreach \x/\lab in {0/jugend, 1.55/funksj., 3.1/postm.}{
  \draw[rounded corners=2pt, fill=black!5] (\x,1.7) rectangle (\x+1.35,2.6);
  \node at (\x+0.675,2.15) {\lab};
}
\node[draw, rounded corners=2pt, anchor=north, text width=42mm, align=center] at (2.2,1.2)
  {\textsc{klassifikator}\\ \scriptsize konvolusjonelt nettverk};
\draw[-{Latex[length=2mm]}] (2.2,1.7) -- (2.2,1.25);
\node[text width=46mm, align=center, font=\scriptsize] at (2.2,0.1)
  {Eit løyst problem: maskina les etiketten på brøkdelen av eit sekund.};

% Skilje
\draw[dashed] (5.1,-0.4) -- (5.1,3.1);

% ---- Høgre: formrom (morfometri) ----
\node[anchor=north west, font=\small\scshape] at (5.6,3.3) {form: eit målbart rom};
\draw[-{Latex[length=2mm]}] (6.0,0.2) -- (6.0,2.9) node[above, font=\scriptsize, align=center] {ornament-\\tettleik};
\draw[-{Latex[length=2mm]}] (6.0,0.2) -- (9.6,0.2) node[right, font=\scriptsize] {setehøgd};
\foreach \px/\py in {6.4/0.7, 6.8/1.5, 7.1/0.9, 7.6/2.2, 7.9/1.1, 8.3/1.8, 8.7/0.6, 9.0/1.4, 8.1/0.5, 7.3/1.9}{
  \fill (\px,\py) circle (1.6pt);
}
\node[text width=40mm, align=center, font=\scriptsize] at (7.8,-0.95)
  {Kontinuerlege, målbare aksar: avstand, retning, variasjon. Metoden faget aldri tok.};
\end{tikzpicture}
\caption{\textsc{Stil mot form.} Stilklassifisering er mønstergjenkjenning på overflatestatistikk — no eit løyst maskinlæringsproblem. Å \emph{måle} form i eit formrom, slik morfometrien gjer, er den kvantitative metoden designhistoria valde bort til fordel for katalogen.}
\end{figure}
